\begin{scanpage}{0207}{189}
\[
\hat E_{0z}=-E_{0z}\frac{\displaystyle\ln\frac{i}{\gamma ka\cos\dfrac{\vartheta_0}2}}{\displaystyle\ln\frac{2i}{\gamma ka\sin\vartheta_0}}.\tag{30.15}
\]
Consequently, for the excitation in an infinite conductor (with $z>z_1$) of current waves (30.06), it is necessary to apply the external field
\[
\begin{gathered}
E_z^e=E_{0z}\frac{\displaystyle\ln\frac{i}{\gamma ka\cos\dfrac{\vartheta_0}2}}{\displaystyle\ln\frac{2i}{\gamma ka\sin\vartheta_0}}
\left[\frac{e^{iw_0z_1}}{2ik\sin^2\dfrac{\vartheta_0}2}\delta(z-z_1)-e^{iw_0z}\varepsilon(z-z_1)\right],\\
\varepsilon(z-z_1)=\begin{cases}1&\text{with }z<z_1,\\0&\text{with }z>z_1.\end{cases}
\end{gathered}\tag{30.16}
\]
In a completely similar way, one may show that the external field
\[
\begin{gathered}
E_z^e=E_{0z}\frac{\displaystyle\ln\frac{i}{\gamma ka\sin\dfrac{\vartheta_0}2}}{\displaystyle\ln\frac{2i}{\gamma ka\sin\vartheta_0}}
\left[\frac{e^{iw_0z_2}}{2ik\cos^2\dfrac{\vartheta_0}2}\delta(z_2-z)-e^{iw_0z}\varepsilon(z_2-z)\right],\\
\varepsilon(z_2-z)=\begin{cases}1&\text{with }z>z_2,\\0&\text{with }z<z_2\end{cases}
\end{gathered}\tag{30.17}
\]
excites the following current in an infinite single-wire line (with $z<z_2$)
\[-Se^{iw_0z_2}\psi_+^0(z_2-z)e^{ik(z_2-z)}.\tag{30.18}\]
Now let us study the diffraction of current waves (30.06) by the semi-infinite conductor $(-\infty,z_2)$. For this purpose, let us use the Lorentz lemma [4]
\[\int\bigl(\mathbf j_1^e\mathbf E_2+\mathbf j_2^m\mathbf H_1\bigr)\,dV=0.\tag{30.19}\]
Here $\mathbf j_1^e=-i\omega\mathbf p_1\delta(\mathbf R-\mathbf R')$ is the current of the auxiliary dipole with the moment $\mathbf p_1$ which is located at point 1 with the coordinates $(R,\vartheta)$; $\mathbf H_1$
\end{scanpage}
